\section{The iPRA framework}
\label{sec:ipra-framework}

\subsection{Why an arithmetic theory can discuss fields and Galois groups}

The only primitive objects of iPRA are natural numbers.  Algebraic objects are
therefore represented by finite codes.  This is not a limitation for the
present proof.  Every polynomial, rational function, permutation, group table,
radical tower, field expression, and derivation appearing below is finite.
The semantic notation is a readable abbreviation for operations on such
codes.

The ProveIt presentation of iPRA has primitive-recursive function symbols and
quantifier-free induction.  Its native scalar calculus is bi-interpretable
with the ordinary first-order presentation; the numerical domain remains all
of $\N$ and equality remains actual equality.  Consequently it is enough to
show the following two facts.

\begin{enumerate}
\item Each construction used in the proof is primitive recursive on codes.
\item Each induction can be stated as an equation or a Boolean invariant about
  a primitive-recursive computation.
\end{enumerate}

The proof does not need induction on an arbitrary semantic property of fields.
It needs induction on lengths of lists, degrees of sparse polynomials, sizes of
finite groups, and lengths of derivations.  All these invariants are decidable.

\subsection{A concrete coding layer}

One convenient coding, not the only possible one, is as follows.

\begin{longtable}{p{0.22\textwidth}p{0.70\textwidth}}
\toprule
Object & Natural-number representation and primitive-recursive operations\\
\midrule
\endhead
Finite sequence & Iterated pairing with stored length.  Entry, append, map,
fold, bounded search, and permutation are primitive recursive.\\
Integer & Zig-zag code: $0,1,-1,2,-2,\ldots$.  Addition and multiplication are
primitive recursive.\\
Rational number & A signed numerator and positive denominator, optionally
normalized by Euclid's algorithm.  Equality is cross multiplication.\\
Monomial & A coefficient and a finite exponent vector.\\
Sparse polynomial & A sorted list of nonzero monomials.  Normalization combines
equal exponent vectors; addition, multiplication, substitution, and formal
derivative are primitive recursive.\\
Rational function & A pair $(P,Q)$ with $Q\ne0$.  Equality of $(P,Q)$ and
$(R,S)$ is the polynomial equation $PS=RQ$.\\
Permutation of $d$ letters & A length-$d$ list containing each number below
$d$ exactly once.  Composition, inverse, sign, cycle decomposition, and
bounded enumeration of all $d!$ permutations are primitive recursive.\\
Finite group & A finite carrier list with multiplication, inverse, and identity
tables, or a subgroup of a permutation group represented by a Boolean
membership table.\\
Field expression & A syntax tree containing constants, generators, $+,-,\cdot$,
inversion with a nonzero side certificate, and named algebraic generators.\\
Radical step & An exponent $m\ge2$, a radicand expression $a$, a new generator
$z$, and the relation $z^m=a$.\\
Derivation & A list of proof nodes in a fixed finite calculus; every node names
its premises and has a local primitive-recursive checker.\\
\bottomrule
\end{longtable}

There is no need to decide all true identities in an arbitrary finitely
presented field.  A certificate carries an explicit derivation.  The checker
only verifies local rules: normalized ring identities, substitution of a
named relation, use of an explicit inverse equation, finite linear algebra,
and finite group computations.  Soundness of the checker is proved by
quantifier-free induction on the derivation length.

\subsection{Proof-carrying field presentations}

A \emph{proof-carrying presentation} consists of:

\begin{itemize}
\item a finite list of generators;
\item a finite list of polynomial and inverse relations;
\item syntax trees for every element used in the argument;
\item a derivation for every asserted equality or nonzero condition;
\item for a finite extension, a finite basis and coordinate calculations for
  multiplication, embeddings, and automorphisms used by the proof.
\end{itemize}

The checker is primitive recursive because it never searches for a proof.
It checks the supplied proof.  This is the same design principle used by a
small proof assistant: semantic truth may be difficult, but checking a finite
proof object is mechanical.

For the generic quintic the base field has an especially simple direct
representation.  Start with sparse polynomials in five variables over $\Q$,
then represent $L=\Q(x_1,\ldots,x_5)$ by fractions.  Equality in $L$ is a
single cross-multiplied polynomial identity.  The subfield
$K=\Q(e_1,\ldots,e_5)$ is represented by rational expressions in five named
symmetric generators together with their embeddings in $L$.  All operations
used in the proof are therefore explicit normalization computations.

\subsection{Radical programs as relational data}

A global operation $z\mapsto\sqrt[m]{z}$ is inappropriate constructively.
Instead we use a finite sequential program with a nondeterministic but
proof-carrying radical instruction:
\[
  z_i^m = a_i(z_0,\ldots,z_{i-1}).
\]
A \emph{run} is a tuple of values in an extension field satisfying every
instruction.  Arithmetic instructions are deterministic; radical
instructions supply a value and a proof of its power equation.  A run is
\emph{complete for $f$} if the designated outputs give the exact factorization
of $f$.

Every ordinary nested-radical expression compiles to such a program.  A case
distinction compiles to a finite decision tree whose guards are decidable
equalities of normalized expressions.  Conversely, a run generates a radical
tower by adjoining its radical values in program order.  Thus the relational
language is neither weaker nor stronger than the standard field-theoretic
notion needed for Abel--Ruffini; it merely exposes the witnesses.

\subsection{Arithmetic predicates used in the final theorem}

The formalization can use the following coded predicates.

\begin{align*}
 \operatorname{RatFunField}(k) &:\quad k\text{ codes the base rational-function presentation},\\
 \operatorname{RadTower}(k,t) &:\quad t\text{ codes a valid radical tower over }k,\\
 \operatorname{Splits}(k,p,t) &:\quad p\text{ factors completely in the top of }t,\\
 \operatorname{Formula}(n,c) &:\quad c\text{ is a degree-}n\text{ radical-program code},\\
 \operatorname{Run}(k,c,w) &:\quad w\text{ is a valid run of }c\text{ over }k.
\end{align*}

The local checking relations can be made primitive recursive.  Their semantic
soundness is a theorem of iPRA about the checker.  The negative theorem is
conveniently factored into two statements:

\begin{align}
 &\forall t\;\neg\operatorname{Splits}(K,P,t),
 \label{eq:no-tower}\\
 &\operatorname{Formula}(5,c)\land\operatorname{Universal}(c)
 \;\longrightarrow\;
   \exists t\;\operatorname{Splits}(K,P,t).
 \label{eq:compile-formula}
\end{align}

Equation~\eqref{eq:no-tower} is the Galois obstruction.
Equation~\eqref{eq:compile-formula} is the primitive-recursive compiler that
specializes a purported coefficient formula to the generic coefficients.
Together they imply that no universal code $c$ exists.

For degrees at most four, fixed constructors produce a radical program and a
factorization derivation from the supplied witness equations.  The internal
statements therefore have the form
\[
 \forall a_0,\ldots,a_n,w\;
   \bigl(\operatorname{WitnessEquations}_n(a,w)=1\bigr)
   \to
   \bigl(\operatorname{FactorCheck}_n(a,w)=1\bigr).
\]
Both sides are equations between primitive-recursive terms, exactly the kind
of induction-free or quantifier-free-induction statement natural to PRA.

\subsection{Where intuitionistic logic matters}

Three classical shortcuts are explicitly removed.

\paragraph{No arbitrary branch choice.}
The square or cube root is a witness, not the result of a choice function.
Cardano's two cube radicals are linked by a product equation.

\paragraph{No unrestricted zero/nonzero split.}
The human proof sometimes writes ``if $s=0$''.  Internally $s$ is a normalized
coded expression in a decidable field presentation, so the split is a Boolean
calculation.  When no decidable equality is available, the primary theorem is
kept as a factorization identity and makes no disjunctive claim.

\paragraph{No appeal to an infinite set of all automorphisms.}
The only automorphism groups used are finite.  For the generic quintic every
automorphism is represented by one of the $120$ permutations of five
symbols.  Radical normal-closure groups are represented by finite tables whose
sizes are bounded by products and factorials extracted from the tower code.

\subsection{Quantifier-free induction is enough}

Examples of the induction statements actually needed are:

\begin{itemize}
\item after processing the first $j$ monomials, the accumulator equals their
  normalized sum;
\item after processing the first $j$ proof nodes, each checked node is sound;
\item after processing the first $j$ radical steps, the generated group-series
  certificate passes its finite checker;
\item after $j$ closure rounds, the Boolean table equals the subgroup generated
  by commutators found so far.
\end{itemize}

Each invariant is an equation between values of primitive-recursive functions.
No induction on a formula containing an unbounded existential quantifier is
required.  Existential mathematical statements are represented by explicit
witness codes produced before the checker is invoked.
